\chapter{Исследовательская часть}

В данном разделе будет произведена постановка задач на исследования, описана методология проведения замеров, приведены и проанализированы их результаты по трём направлениям: исследование оптимальной конфигурации эволюционных операторов, исследование сходимости комбинированного метода относительно муравьиного алгоритма и исследование  устойчивости к выбору параметров $\alpha$ и $\beta$, а также относительно муравьиного алгоритма.

\section{Исследование вариантов метода}

Объектом данного исследования является разработанный муравьиный алгоритм с динамическими гетерогенными параметрами на основе генетического алгоритма.

Предметом данного исследования является зависимость качества получаемых алгоритмом решений в зависимости от конфигурации эволюционных операторов.

\textbf{Цель исследования} -- определить оптимальную конфигурацию эволюционных операторов, обеспечивающую наилучшее качество решений задач коммивояжёра разработанным алгоритмом.

\subsection{Описание исследования}

Варьируемыми факторами исследования являются:
\begin{itemize}
	\item стратегия селекции -- элитарная, турнирная, рулеточная;
 	\item стратегия скрещивания -- арифметическая, BLX, SBX;
	\item стратегия мутации -- равномерная, гауссовская.
\end{itemize}

Таким образом требуется сравнить $3 * 3 * 2 = 18$ конфигураций разработанного алгоритма.

Также варьируется входной взвешенный неорграф, описывающий задачу коммивояжёра. Были выбраны 6 тестовых задач из библиотеки TSPLIB: \texttt{ulysses22.tsp}, \texttt{swiss42.tsp}, \texttt{eil51.tsp}, \texttt{berlin52.tsp}, \texttt{st70.tsp}, \texttt{eil76.tsp}. Данные задачи различаются размерностью (от 22 до 76 городов), топологией и методом задания весов, что позволяет оценить алгоритм на различных типах входных данных.

Постоянными параметрами и их значениями в данном исследовании являются: 
\begin{itemize}
	\item начальные значения $\alpha$ и $\beta$: 1 и 3 соответственно;
	\item коэффициент испарения $p$: 0.5;
	\item множитель феромона $Q$: 1;
	\item количество муравьёв $m$: $n$
	\item начальный уровень феромона $\tau_0$: $\frac{m}{C_{nn}}$
	\item количество поколений: 300;
	\item период адаптации $w$: 5;
	\item количество родителей: $0.4m$.
\end{itemize}
Для параметров эволюционных стратегий были выбраны следующие значения:
\begin{itemize}
	\item равномерная мутация: от -1 до 1;
	\item гауссова мутация: $\sigma=0.05$, $\mu=0$;
	\item турнирная селекция: $k=5$;
	\item BLX-скрещивание: $\gamma=0.5$;
	\item SBX-скрещивание: $\eta=2$.
\end{itemize}

Для сравнения конфигураций предлагаются следующие метрики:

\textbf{Статистические оценки ошибки}

Относительная ошибка $\varepsilon$~(\ref{eq:1var}) -- отношение разности длины найденного маршрута с лучшим известным решением для данного графа к этому же лучшему решению.
\begin{equation}
	\varepsilon = \frac{|L_{\text{найденное}} - L_{\text{лучшее}}|}{L_{\text{лучшее}}}.
	\label{eq:1var}
\end{equation}

Средняя относительная ошибка $\mu_{err}$~(\ref{eq:1muerr}) -- среднее арифметическое относительной ошибки по всем графам и запускам:
\begin{equation}
\mu_{err} = \frac{1}{R*(|G|-1)} \sum_{r=1}^{R}\sum_{g}^{G} \varepsilon_{rg},
\label{eq:1muerr}
\end{equation}
где
\begin{itemize}
	\item $R$ -- количество запусков для заданной конфигурации и графа;
	\item $G$ -- множество входных графов.
\end{itemize}

Среднеквадратичное отклонение относительной ошибки $\sigma_{err}$~(\ref{eq:1sigmaerr}) -- показатель стабильности алгоритма.
\begin{equation}
\sigma_{err} = \sqrt{\frac{1}{R*(|G|-1)} \sum_{r=1}^{R}\sum_{g}^{G} (\varepsilon_{rg} - \mu_{err})^2}.
\label{eq:1sigmaerr}
\end{equation}

Использование относительной ошибки позволяет нивелировать различные размерности графов и оценивать их единым образом. $\mu_{err}$ и $\sigma_{err}$ фактически являются оценками мат. ожидания и стандартного отклонения относительной ошибки метода в заданной конфигурации.

\textbf{Ранги конфигураций}

Как метрику для сравнения конфигураций между собой предлагается использовать средний ранг. Для каждого графа $g$ все конфигурации упорядочиваются по возрастанию длины найденного ими решения. Лучшей конфигурации на данном графе присваивается ранг $r_g = 1$, следующей -- ранг 2 и так далее до 18.

Средний ранг конфигурации вычисляется как среднее арифметическое её рангов по всем графам~(\ref{eq:1rank}).
\begin{equation}
\bar{r} = \frac{1}{|G|} \sum_{g=1}^{G} r_g.
\label{eq:1rank}
\end{equation}

Ранжирование позволит оценить качество решений полученных конфигураций устойчиво относительно масштаба и выбросов: на каждой задаче конфигурации сравниваются только между собой, что важно при объединении результатов с различных графов, имеющих разную размерность и структуру.

\textbf{Алгоритм исследования}

\begin{enumerate}
	\item Для каждой тестовой задачи:
	\begin{enumerate}
		\item загрузить граф из TSP-файла;
		\item вычислить жадный тур для инициализации феромонов;
	\end{enumerate}
	\item Для каждой из 18 конфигураций:
	\begin{enumerate}
		\item выполнить $R$ независимых запусков алгоритма;
		\item для каждого запуска зафиксировать лучшее найденное решение (score);
	\end{enumerate}
	\item Для каждой задачи и каждой конфигурации:
	\begin{enumerate}
		\item вычислить относительную ошибку $\varepsilon_{rg}$;
		\item рассчитать среднее $\mu_{err}$ и стандартное отклонение $\sigma_{err}$ по $R$ запускам;
		\item рассчитать ранг конфигурации по задаче;
	\end{enumerate}
	\item Для каждой конфигурации:
	\begin{enumerate}
		\item рассчитать средний ранг по всем задачам (ранг 1 присваивается лучшему методу на задаче);
	\end{enumerate}
\end{enumerate}

\subsection{Результаты исследования}

Для каждой конфигурации эволюционных операторов и входного графа алгоритм запускался $R = 20$ раз. Усреднённые результаты представлены в таблицах~\ref{tbl:berlin52.tsp}-\ref{tbl:ulysses22.tsp}, а результаты ранжирования и расчёта относительных ошибок представлены в таблице~\ref{tbl:algorithm_ranking}. На рисунке~\ref{fig:variants-bar} представлена столбчатая диаграмма средних рангов конфигураций, а на рисунке~\ref{fig:variants-heat} -- тепловая карта рангов конфигураций по входным графам.

\begin{figure}[H]
	\centering
	\includegraphics[width=\textwidth]{variants/rank}
	\caption{Диаграмма средних рангов конфигураций эволюционных операторов}
	\label{fig:variants-bar}
\end{figure}

\begin{figure}[H]
	\centering
	\includegraphics[width=0.8\textwidth]{variants/heatmap}
	\caption{Тепловая карта рангов конфигураций эволюционных операторов по входным графам}
	\label{fig:variants-heat}
\end{figure}

Таким образом как по среднему рангу, так и по средней относительной ошибке, наиболее удачной конфигурацией является сочетание рулеточной селекции, BLX-скрещивания и гауссовской мутации. Данный вариант будет использован в дальнейших исследованиях как основной.

\section{Исследование сходимости метода}

Объектом данного исследования является разработанный гетерогенный муравьиный алгоритм с динамическими гетерогенными параметрами на основе генетического алгоритма с оптимальными эволюционными операторами, определёнными в предыдущей секции.

Предметом данного исследования является сходимость алгоритма к субоптимальным решениям в зависимости от количества поколений и наличия локальной оптимизации.

\textbf{Цель исследования} -- оценить влияние количества поколений и локальной оптимизации на сходимость разработанного алгоритма, а также сравнить его сходимость со сходимостью классического муравьиного алгоритма.

\subsection{Описание исследования}

Варьируемыми факторами исследования являются:

\begin{itemize}
	\item \textbf{Количество поколений}: 500, 800, 1000;
	\item \textbf{Тип алгоритма}:
	\begin{itemize}
		\item классический муравьиный алгоритм;
		\item разработанный алгоритм без локальной оптимизации;
		\item разработанный алгоритм с 2-opt локальной оптимизацией.
	\end{itemize}
\end{itemize}

Также варьируется входной взвешенный неорграф, описывающий задачу коммивояжёра. Были выбраны 7 тестовых задач из библиотеки TSPLIB: \texttt{swiss42.tsp}, \texttt{eil51.tsp}, \texttt{berlin52.tsp}, \texttt{eil76.tsp}, \texttt{rd100.tsp}, \texttt{eil101.tsp}, \texttt{tsp225.tsp}. Прежде всего конфигурации существенно отличаются размерностью графа (от 42 до 225), что позволит оценить эффективность алгоритма при различных размерах входных данных.

Постоянными параметрами и их значениями в данном исследовании являются: 
\begin{itemize}
	\item начальные значения $\alpha$ и $\beta$ (для классического алгоритма постоянные): 1 и 5 соответственно;
	\item коэффициент испарения $p$: 0.5;
	\item множитель феромона $Q$: 1;
	\item количество муравьёв $m$: $n$
	\item начальный уровень феромона $\tau_0$: $\frac{m}{C_{nn}}$
	\item период адаптации $w$: 5;
	\item количество родителей: $0.4m$;
	\item эволюционные стратегии: рулеточная селекция, BLX-кроссинговер, гауссовская мутация.
\end{itemize}

Для оценки сходимости предлагаются следующие метрики:
\begin{itemize}
	\item количество итераций до лучшего решения (ITB) -- номер итерации, на которой алгоритм впервые достигает наилучшего найденного им решения. Меньшее значение ITB свидетельствует о более быстрой сходимости;
	\item длина наилучшего маршрута -- длина наилучшего маршрута, найденного алгоритмом за все поколения;
	\item время затраченное методом.
\end{itemize}

\textbf{Алгоритм исследования}

\begin{enumerate}
	\item Для каждой тестовой задачи:
	\begin{enumerate}
		\item загрузить граф из TSP-файла;
		\item вычислить жадный тур для инициализации феромонов.
	\end{enumerate}
	\item Для каждого количества поколений:
	\begin{enumerate}
		\item Для каждого типа алгоритма:
		\begin{enumerate}
			\item выполнить $R$ независимых запусков;
			\item для каждого запуска зафиксировать:
			\begin{itemize}
				\item длину найденного маршрута;
				\item количество итераций до лучшего решения (ITB);
				\item время затраченное алгоритмом;
			\end{itemize}
		\end{enumerate}
	\end{enumerate}
	\item Для каждого количества поколений и типа алгоритма вычислить средние и стандартные отклонения метрик по 20 запускам.
\end{enumerate}

\subsection{Результаты исследования}

Исследование проводились на устройстве со следующими характеристиками:

\begin{itemize}
	\item Процессор: AMD Ryzen 7 5800H (16 ядер, 4.46 ГГц);
	\item Оперативная память: 16 ГБ;
	\item Операционная система: openSUSE Tumbleweed x86\_64 (Linux 6.19.12-1-default).
\end{itemize}

Для каждой пары алгоритма и входного графа алгоритм запускался $R = 20$ раз. Усреднённые результаты запусков представлены в таблице~\ref{tbl:convergence}. На рисунке~\ref{fig:convergence-boxplot} представлена коробчатая диаграмма, отображающая распределение медианы ITB полученным алгоритмами в задачах, а на рисунке~\ref{fig:convergence-scatter} -- точечный график зависимости медиан длины маршрута и ITB разными алгоритмами по задачам.

\begin{figure}[H]
	\centering
	\includegraphics[width=\textwidth]{convergence/boxplots}
	\caption{Коробчатая диаграмма распределения ITB по задачам и алгоритмам}
	\label{fig:convergence-boxplot}
\end{figure}

\begin{figure}[H]
	\centering
	\includegraphics[width=\textwidth]{convergence/scatter}
	\caption{Точечный график зависимости ITB и длины найденного маршрута по задачам и алгоритмам}
	\label{fig:convergence-scatter}
\end{figure}

Разработанный алгоритм без локальной оптимизации (мод.) демонстрирует качество, хоть и несколько худшее, но сопоставимое с классическим муравьиным алгоритмом (разница в пределах 1–2\%). Однако достигает этого качества меньшим количеством итераций: на eil76 ускорение по ITB составило 3.7 раза (49.5 против 183.0 итераций), на eil101 -- до 8 раз (39.0 против 311.0), на tsp225 -- 5.2 раза (45.5 против 235.0). С точки зрения быстродействия, модификация просчитывает тоже число поколений за время приблизительно на 10\% большее, чем классический метод.

Добавление 2-opt локальной оптимизации (мод. + лок.) позволяет значительно улучшить качество решений: на berlin52 достигнут глобальный оптимум 7542.00, на eil101 улучшение составило 7.8\% относительно классического ACO (637.90 против 691.50), на tsp225 -- 6.2\% (3994.10 против 4259.60). Однако это достигается ценой увеличения времени выполнения в 1.3-1.7 раза. При этом показатель ITB также выше в версии с локальной оптимизацией, за счёт того, что эволюционные операторы позволяют дольше покрывать большое пространство решений, а локальная оптимизация -- находить оптимальные решения.

Таким образом, разработанный алгоритм без локальной оптимизации обеспечивает более быструю сходимость по сравнению с классическим муравьиным алгоритмом, ценой не существенного уменьшения качества получаемых решений. Добавление 2-opt локальной оптимизации существенно улучшает решение, однако увеличивает вычислительные затраты.

\section{Исследование устойчивости метода}

Объектом данного исследования является разработанный гетерогенный муравьиный алгоритм с динамическими гетерогенными параметрами на основе генетического алгоритма с оптимальными эволюционными операторами, определёнными ранее.

Предметом данного исследования является зависимость качества решений, получаемых алгоритмом, от начальных значений параметров $\alpha$ и $\beta$, управляющих влиянием феромона и эвристической информации.

\textbf{Цель исследования} -- оценить устойчивость разработанного алгоритма к выбору параметров $\alpha$ и $\beta$ в сравнении с классическим муравьиным алгоритмом.

\subsection{Описание исследования}

Варьируемыми факторами исследования являются:

\begin{itemize}
	\item начальное значение параметра $\alpha$: от 0.2 до 2.0 с шагом 0.2;
	\item начальное значение параметра $\beta$: от 0.5 до 5.0 с шагом 0.5;
	\item тип алгоритма: классический муравьиный алгоритм и разработанная модификация.
\end{itemize}

Диапазоны выбирались вокруг рекомендуемых значений из~\cite{dorigo-ant-book}, выходя за их пределы для оценки качества решений при не эффективных параметрах.

В качестве входного графа был выбран файл в формате TSPLIB eil76.tsp, имеющей размерность в 76 узлов, что более существенно оценить зависимость от параметров.

Постоянными параметрами и их значениями в данном исследовании являются: 
\begin{itemize}
	\item коэффициент испарения $p$: 0.5;
	\item множитель феромона $Q$: 1;
	\item количество муравьёв $m$: $n$
	\item начальный уровень феромона $\tau_0$: $\frac{m}{C_{nn}}$
	\item количество поколений: 500;
	\item период адаптации $w$: 5;
	\item количество родителей: $0.4m$;
	\item эволюционные стратегии: рулеточная селекция, BLX-кроссинговер, гауссовская мутация;
\end{itemize}
где $n$ -- количество вершин в графе, $C_{nn}$ -- длина маршрута, построенная жадным алгоритмом.


Для оценки устойчивости метода предлагаются следующие 2 метрики:
$r_{\mu_{L}}$~(\ref{eq:3r}) -- разность между максимальным и минимальным значениями средней длины маршрута.:
\begin{equation}
r_{\mu_{L}} = \max \mu_{L} - \min \mu_{L}.
\label{eq:3r}
\end{equation}

$CV_{\mu_{L}}$~(\ref{eq:3cv}) -- относительное среднеквадратичное отклонение, отношение стандартного отклонения средних длин маршрута к их среднему значению по всем комбинациям параметров.
\begin{equation}
CV_{\mu_{L}} = \frac{\sigma_{\mu_{L}}}{\mu_{\mu_{L}}},
\label{eq:3cv}
\end{equation}
где
\begin{itemize}
	\item $\mu_{L}$~(\ref{eq:3mu}) -- средняя длина маршрута по $R$ запускам для заданной комбинации $(\alpha, \beta)$;
	\begin{equation}
		\mu_{L}(\alpha, \beta) = \frac{1}{R} \sum_{r=1}^{R} L_r(\alpha, \beta).
		\label{eq:3mu}
	\end{equation}
	
	\item $\sigma_{L}$~(\ref{eq:3sigma}) -- среднеквадратичное отклонение длины маршрута по $R$ запускам для заданной комбинации $(\alpha, \beta)$.
	\begin{equation}
		\sigma_{L}(\alpha, \beta) = \sqrt{\frac{1}{R-1} \sum_{r=1}^{R} (L_r - \mu_{L})^2}.
		\label{eq:3sigma}
	\end{equation}
\end{itemize}

В данном случае $CV$ предлагается как интегральная оценка стабильности алгоритма по всем парам $\alpha$ и $\beta$ в совокупности. Чем меньше $CV$ тем менее разбросаны средние значения путей относительно друг друга и соответственно алгоритм более стабильный.

\textbf{Алгоритм исследования}

\begin{enumerate}
	\item Загрузить граф \texttt{eil76.tsp} и вычислить жадный тур для инициализации феромонов.
	\item Для каждого значения $\alpha$:
	\begin{enumerate}
		\item Для каждого значения $\beta$:
		\begin{enumerate}
			\item Для каждого типа алгоритма:
			\begin{itemize}
				\item выполнить $R$ независимых запусков;
				\item вычислить $\mu_{L}(\alpha, \beta)$ и $\sigma_{L}(\alpha, \beta)$
			\end{itemize}
		\end{enumerate}
	\end{enumerate}
	\item Для каждого типа алгоритма вычислить $CV_{\mu_{L}}$ и $r_{\mu_{L}}$.
\end{enumerate}


\subsection{Результаты исследования}

Для каждой пары параметров $\alpha$ и $\beta$ оба алгоритма запускался $R = 20$ раз. Усредненные результаты запусков представлены в таблице~\ref{tbl:stability}. На рисунке~\ref{fig:stability-diff} изображены графики зависимости средних длин маршрутов полученных муравьиным алгоритмом и разработанной модификацией от параметров $\alpha$ и $\beta$.

\begin{figure}[H]
	\centering
	\includegraphics[width=\textwidth]{stability/comparison}
	\caption{Сравнение средних длин маршрутов, полученных муравьиным алгоритмом и разработанной модификацией при варьировании $\alpha$ и $\beta$}
	\label{fig:stability-diff}
\end{figure}

\textbf{Муравьиный алгоритм}

Размах $\mu_{L}$: 1101.85

Коэффициент вариации $\mu_{L}$: 27.96\%

\textbf{Разработанная модификация}

Размах $\mu_{L}$: 124.95

Коэффициент вариации $\mu_{L}$: 4.54\%

Разработанная модификация демонстрирует значительно более высокую устойчивость к выбору параметров $\alpha$ и $\beta$ по сравнению с классическим муравьиным алгоритмом. Коэффициент вариации средней длины маршрута снижен с 27.96\% до 4.54\%, а размах -- со 1101.85 до 124.95 единиц. Это означает, что предложенный алгоритм сохраняет высокое качество решений независимо от параметров, а классическая модификация требует более тонкой настройки.

\section*{Вывод}

В рамках исследовательского раздела было проведено 3 исследования над разработанной модификацией муравьиного алгоритма.

В первом исследовании проведён подбор оптимальной конфигурации эволюционных операторов для разработанного гетерогенного муравьиного алгоритма. На основе анализа среднего ранга и относительной ошибки по 6 тестовым задачам TSPLIB лучшей признана конфигурация: рулеточная селекция, BLX-кроссовер и гауссовская мутация.

Во втором исследовании проведена оценка сходимости разработанного алгоритма в сравнении с классическим муравьиным алгоритмом. Разработанный алгоритм без локальной оптимизации показал сопоставимое качество решений (разница 1–2\%), однако достигает его в 3.7–8 раз быстрее по числу итераций. Добавление 2-opt локальной оптимизации улучшает качество решений на 6–8\%, но ценой увеличения времени выполнения в 1.3–1.7 раза.

В третьем исследовании проведена оценка устойчивости разработанного алгоритма к выбору параметров $\alpha$ и $\beta$ на графе eil76. Разработанная модификация показала значительное превосходство: коэффициент вариации средней длины маршрута снижен с 27.96\% до 4.54\%, а размах -- с 1101.85 до 124.95 единиц, что делает алгоритм пригодным для практического использования в условиях, когда оптимальные параметры задачи заранее неизвестны.

\clearpage
